%!TEX root = main.tex

\begin{abstract}

Simulation offers unique values for both enumeration and extrapolation purposes, and is becoming increasingly important for managing the massive machine learning (ML) clusters and large-scale distributed training jobs.
In this paper, we build \sys to tackle three key challenges in large-scale training simulation:
(1) tracing the runtime training workloads at each device in an ex-situ fashion so we can use a single device to obtain the actual execution graphs of 1K-GPU training, (2) accurately estimating the collective communication without high overheads of discrete-event based network simulation, and (3) accounting for the interference-induced computation slowdown from overlapping communication and computation kernels on the same device.
\sys supports both dense and MoE workloads with high fidelity and further enables GPU memory–aware ex-situ simulation. 
\sys delivers on average 8\% error in training step---$\sim$3x lower than state-of-the-art simulators---for GPT-175B on a 96-GPU H800 cluster with 3D parallelism on Megatron-LM under 2 minutes.
% On a 96-GPU H800 cluster, \sys achieves an average step-time error of 8\% for GPT-175B with 3D parallelism on Megatron-LM, delivering $\sim$3$\times$ lower error than state-of-the-art simulators, all within 2 minutes.  
% It can simulate a 8,192-GPU setting in 1.38 hours compared to 5.2 days for recent work using ns-3. 
% The scaling law have made it standard practice to train large-scale deep neural network (DNN) models using thousands of GPUs, leveraging state-of-the-art techniques. This approach necessitates a precise simulator to assist with the design and optimization of the training process before actual model training begins, significantly reducing both time and financial costs throughout the development pipeline. However, existing simulators fail to meet these demands due to their inability to accurately model computation and communication workloads in large-scale training scenarios.
% In this paper, we re-examine existing simulation approaches and, based on this analysis, introduce our simulator, Moye. To closely replicate the training process, we delve into mainstream training frameworks and communication libraries, developing a suite of tools that facilitate accurate construction of training workloads, prediction of communication times, and simulation of performance slowdowns caused by kernel overlap. We adopt a profile-and-predict approach to derive the iteration time of large-scale distributed tasks from knowledge obtained from small-scale tasks. We evaluate the performance of the simulator on NVIDIA H800 and A800 clusters, demonstrating an average prediction accuracy of up to 92.25\%.




\end{abstract}
